Let $\varphi$ be Euler's totient function.
For integers $m\ge0$ and $k\ge1$, consider the length-$k$ block
$$(\varphi(m+1),\varphi(m+2),\dots,\varphi(m+k)).$$
An \emph{ordering pattern} of length $k$ means a permutation $\pi\in S_k$ giving the relative order of these $k$ numbers by size.
Because the definition in the problem file says ``the $k!$ possible ordering patterns'', the natural strict interpretation is:
\begin{quote}
we only count a block if the $k$ totients are pairwise distinct, and then the pattern is the unique permutation that sorts them increasingly.
\end{quote}
Define $F(n)$ as the largest $k$ such that \emph{every} permutation in $S_k$ occurs as the strict ordering pattern of some block with $m+k\le n$.

Questions asked in the problem file:
\begin{enumerate}
  \item Is $F(n)=(c+o(1))\log\log\log n$ for some constant $c$?
  \item Is the first pattern which fails to appear always the monotone decreasing pattern
  $$\varphi(m+1)>\varphi(m+2)>\cdots>\varphi(m+k)?$$
  \item If one allows equality in patterns, is the ``natural'' ordering mimicking $\varphi(1),\dots,\varphi(k)$ the most likely to appear?
\end{enumerate}
